\clearpage
\begin{explanationbox}
\textbf{第一步，先认初中几何。}暂时没有抛物线、焦点和坐标；只看普通斜线上的两个直角三角形。
\begin{center}\begin{tikzpicture}[font=\small,>=stealth]
\draw[rounded corners=3pt,gray!45] (0,0) rectangle (14.2,12);
\node at (7.1,11.45) {先认熟悉的结构：两个直角三角形，相同的倾斜程度};\begin{scope}[shift={(1.7,4.2)},scale=1.25]
\fill[blue!7] (1,0)--(4,0)--(4,4)--cycle;
\fill[orange!10] (1,0)--(.25,0)--(.25,-1)--cycle;
\draw[blue!75!black,very thick] (1,0)--(4,4);
\draw[orange!85!black,very thick] (1,0)--(.25,-1);
\draw[gray,dashed] (4,0)--(4,4) node[midway,right=7pt,text=blue!75!black] {$u$};
\draw[gray,dashed] (.25,0)--(.25,-1) ;\draw[gray!60,thin] (.25,0)--(-.45,0) (.25,-1)--(-.45,-1);\draw[<->,orange!85!black] (-.4,0)--(-.4,-1) node[midway,left=6pt] {$v$};
\draw[blue!60!black,thick] (1,0)--(4,0) node[midway,below=7pt] {$h$};
\draw[orange!85!black,thick] (.25,0)--(1,0) node[midway,above=6pt] {$k$};
\draw[gray] (3.8,0)--(3.8,.2)--(4,.2);
\draw[gray] (.43,0)--(.43,-.18)--(.25,-.18);
\draw[violet] (1.36,0) arc[start angle=0,end angle=53.13010235,radius=.36];
\draw[violet] (.64,0) arc[start angle=180,end angle=233.13010235,radius=.36];
\fill (4,4) circle(2pt) node[above=8pt] {$A$};
\fill (.25,-1) circle(2pt) node[below=8pt] {$B$};
\fill (4,0) circle(1.5pt) node[below right=6pt] {$H$};
\fill (.25,0) circle(1.5pt) node[above=16pt] {$K$};

\fill (1,0) circle(2pt) node[below right=6pt] {$P$};
\end{scope}
\node[anchor=west,align=left] at (8.6,7.1)
 {$A,P,B$ 共线\\[12pt]$H,P,K$ 共线\\[12pt]
 $AH\parallel BK$\\[16pt]两处直角相等\\[12pt]$P$ 处两角是对顶角\\[12pt]所以两个三角形相似};
\node at (7.1,1.7) {四条标出的边长 $h,k,u,v$ 都是正数。};
\node at (7.1,.65) {$\displaystyle \triangle PHA\sim\triangle PKB\qquad\Longrightarrow\qquad\frac hu=\frac kv$};
\end{tikzpicture}
\end{center}
对应的水平边与竖直边之比相等：$h/u=k/v$。先记住这一个熟悉的相似关系，四个字母都表示正的边长。
\end{explanationbox}
\clearpage
\begin{explanationbox}
\textbf{第二步，把同一结构放回抛物线。}图形、尺度和颜色保持不变；普通点 $P$ 对应焦点 $F$，两端点落在曲线上。
\begin{center}\begin{tikzpicture}[font=\small,>=stealth]
\draw[rounded corners=3pt,gray!45] (0,0) rectangle (14.2,12);
\node at (7.1,11.45) {再加抛物线：普通点对应焦点，边长对应坐标差};\begin{scope}[shift={(1.7,4.2)},scale=1.25]
\draw[->,gray] (-.45,0)--(4.75,0) node[above] {$x$};
\draw[->,gray] (0,-1.8)--(0,4.75) node[left] {$y$};
\draw[teal!75!black,thick,samples=120,domain=-1.7:4.5,variable=\t] plot ({\t*\t/4},{\t});

\fill[blue!7] (1,0)--(4,0)--(4,4)--cycle;
\fill[orange!10] (1,0)--(.25,0)--(.25,-1)--cycle;
\draw[blue!75!black,very thick] (1,0)--(4,4);
\draw[orange!85!black,very thick] (1,0)--(.25,-1);
\draw[gray,dashed] (4,0)--(4,4) node[midway,right=7pt,text=blue!75!black] {$u$};
\draw[gray,dashed] (.25,0)--(.25,-1) ;\draw[gray!60,thin] (.25,0)--(-.45,0) (.25,-1)--(-.45,-1);\draw[<->,orange!85!black] (-.4,0)--(-.4,-1) node[midway,left=6pt] {$v$};
\draw[blue!60!black,thick] (1,0)--(4,0) node[midway,below=7pt] {$h$};
\draw[orange!85!black,thick] (.25,0)--(1,0) node[midway,above=6pt] {$k$};
\draw[gray] (3.8,0)--(3.8,.2)--(4,.2);
\draw[gray] (.43,0)--(.43,-.18)--(.25,-.18);
\draw[violet] (1.36,0) arc[start angle=0,end angle=53.13010235,radius=.36];
\draw[violet] (.64,0) arc[start angle=180,end angle=233.13010235,radius=.36];
\fill (4,4) circle(2pt) node[above=8pt] {$A$};
\fill (.25,-1) circle(2pt) node[below=8pt] {$B$};
\fill (4,0) circle(1.5pt) node[below right=6pt] {$H$};
\fill (.25,0) circle(1.5pt) ;\node at (.7,1.0) {$K$};\draw[gray,thin] (.58,.78)--(.25,.06);

\fill (1,0) circle(2pt) node[below right=6pt] {$F$};
\node[above left=8pt] at (0,0) {$O$};
\end{scope}
\node[anchor=west,align=left] at (8.5,7.2)
 {此图取 $p=2$\\[12pt]$P\ \longrightarrow\ F(1,0)$\\[14pt]
 $A(4,4),\quad B(1/4,-1)$\\[14pt]
 $u=y_A=4$\\[10pt]$v=-y_B=1$\\[10pt]
 $h=x_A-1=3$\\[10pt]$k=1-x_B=3/4$};
\node at (7.1,1.55) {注意：$v$ 是正的边长，下端点的纵坐标却是 $-v$。};
\node at (7.1,.6) {$\displaystyle \frac hu=\frac kv\quad\text{保留；再用}\quad x=\frac{y^2}{2p}\quad\text{替换横坐标。}$};
\end{tikzpicture}
\end{center}
对一般非竖直焦点弦，负斜率时先上下翻折，再按图中的左右端点标号：$u=y_A>0$、$v=-y_B>0$。横向边长为
\[
 h=\frac{u^2-p^2}{2p},\qquad k=\frac{p^2-v^2}{2p}.
\]
相似给 $h/u=k/v$，整理成 $(u+v)(uv-p^2)=0$。因 $u+v>0$，得到 $uv=p^2$，所以\textbf{带符号的纵坐标积为 $y_Ay_B=-uv=-p^2$}。完整推导见方法二。
\end{explanationbox}
\clearpage
\begin{explanationbox}
\textbf{第三步，观察变化，再抽象成定值。}三图只改变过焦点直线的方向。端点可以不对称；不要把“乘积固定”误读成“两数互为相反数”。
\begin{center}\begin{tikzpicture}[font=\small,>=stealth]
\draw[rounded corners=3pt,gray!45] (0,0) rectangle (14.2,12.7);
\node at (7.1,12.15) {同一抛物线、同一尺度：端点在变，两个坐标积不变};\node at (2.25,11.0) {通径：上下对称};
\begin{scope}[shift={(1.6,6.7)},scale=.65]
\draw[teal!75!black,thick,samples=120,domain=-4.5:4.5,variable=\t] plot ({\t*\t/4},{\t});
\draw[->,gray] (-.7,0)--(5.25,0) node[above] {$x$};
\draw[->,gray] (0,-4.8)--(0,4.9);
\fill (1,0) circle(1.7pt) node[below right=5pt] {$F$};
\draw[blue!75!black,very thick] (1.0,2)--(1,0);\draw[orange!85!black,very thick] (1,0)--(1.0,-2);
\fill (1.0,2) circle(1.7pt) node[above=7pt] {$A$};\fill (1.0,-2) circle(1.7pt) node[below=7pt] {$B$};
\end{scope}
\node at (2.25,2.6) {$y_Ay_B=2\times(-2)=-4$};
\node at (2.25,1.65) {$x_Ax_B=1\times 1=1$};
\node at (6.65,11.0) {斜弦：上长下短};
\begin{scope}[shift={(6.0,6.7)},scale=.65]
\draw[teal!75!black,thick,samples=120,domain=-4.5:4.5,variable=\t] plot ({\t*\t/4},{\t});
\draw[->,gray] (-.7,0)--(5.25,0) node[above] {$x$};
\draw[->,gray] (0,-4.8)--(0,4.9);
\fill (1,0) circle(1.7pt) node[below right=5pt] {$F$};
\draw[blue!75!black,very thick] (4.0,4)--(1,0);\draw[orange!85!black,very thick] (1,0)--(0.25,-1);
\fill (4.0,4) circle(1.7pt) node[above=7pt] {$A$};\fill (0.25,-1) circle(1.7pt) node[below=7pt] {$B$};
\end{scope}
\node at (6.65,2.6) {$y_Ay_B=4\times(-1)=-4$};
\node at (6.65,1.65) {$x_Ax_B=4\times \frac14=1$};
\node at (11.05,11.0) {斜弦：上短下长};
\begin{scope}[shift={(10.4,6.7)},scale=.65]
\draw[teal!75!black,thick,samples=120,domain=-4.5:4.5,variable=\t] plot ({\t*\t/4},{\t});
\draw[->,gray] (-.7,0)--(5.25,0) node[above] {$x$};
\draw[->,gray] (0,-4.8)--(0,4.9);
\fill (1,0) circle(1.7pt) node[above right=7pt] {$F$};
\draw[blue!75!black,very thick] (0.25,1)--(1,0);\draw[orange!85!black,very thick] (1,0)--(4.0,-4);
\fill (0.25,1) circle(1.7pt) node[above=7pt] {$A$};\fill (4.0,-4) circle(1.7pt) node[below=7pt] {$B$};
\end{scope}
\node at (11.05,2.6) {$y_Ay_B=1\times(-4)=-4$};
\node at (11.05,1.65) {$x_Ax_B=\frac14\times 4=1$};
\node at (7.1,.65) {$y^2=4x,\ p=2$：负积说明纵坐标异号，不要求绝对值相等。};
\end{tikzpicture}
\end{center}
读图时，上方纵坐标为正，下方为负；绝对值相乘均为 $4$，所以纵积均为 $-4$。横积则均为 $1$。
一般情形中，只需证明纵积，横积就随之得到：
\[
 x_1x_2=\frac{y_1^2y_2^2}{4p^2}=\frac{(y_1y_2)^2}{4p^2}=\frac{p^2}{4}.
\]
已知一端纵坐标 $y_A=t\ne0$，另一端由乘积配对确定：$y_B=-p^2/t$，$x_B=p^3/(2t^2)$。上面的三图用于观察，下面的证明覆盖任意有限焦点弦。
\end{explanationbox}
